\section{Details on human evaluation}
\label{sec:app-human-eval}

To end-goal of \approach{} is to improve the \textit{attribution} of generation models through post-editing while \textit{preserving} the original intent.
Attribution and preservation are both subjective properties that may change with even small edits.
In the main paper, we present two automatic metrics to conveniently gauge these properties, but rely on a human evaluation as the gold standard.
In this section, we describe how we conducted the human evaluation and what instructions and examples annotators were provided.

\paragraph{Rater recruitment and training}
We engaged with a vendor supplier of full-time crowd workers to recruit human annotators for our task.
Annotators were asked to review the instructions below and were provided direct feedback on their responses during the pilot annotation runs.
We had 3 annotators rate each example in the pilot phase to measure inter-annotator agreement, and had a single rater annotate each example afterwards.

\subsection{Instructions: Overview}

In this task you will evaluate the quality of text generated by a system (\textbf{the ``passage''}) based on how well it represents information from multiple pieces of \textbf{``evidence''}.

We will be using two categories to evaluate the quality of the passage: \textbf{Attribution} and \textbf{Intent Similarity}. You will evaluate these categories in succession. In some tasks, you will only evaluate Attribution. The task interface will guide you through the flow; you can also see the overall task flow in the diagram below. 

\textbf{Note:} The passage may appear very fluent and well-formed, but still contain slight inaccuracies that are not easy to discern at first glance. Pay close attention to the text. Read it carefully as you would when proofreading.

\subsection{Instructions: Attribution}

\newcommand{\attrspan}[1]{\sethlcolor{lime}\hl{#1}}
\newcommand{\nattrspan}[1]{\sethlcolor{pink}\hl{#1}}

\begin{figure*}[!thp]
    \centering
    \includegraphics[width=\textwidth]{figures/human-eval-attribution}
    \caption{
    Screenshot of interface to annotate attribution at the sentence level.
    annotators were asked to mark sentences as being \attrspan{fully attributable} or \nattrspan{not fully attributable} by clicking each sentence, and rating each piece of evidence as being useful or not in helping determine attribution of the passage. Annotators were also presented with the context of the generation.
    }
    \label{fig:human-eval-attribution}
\end{figure*}
\begin{figure*}[!thp]
    \centering
    \includegraphics[width=\textwidth]{figures/human-eval-preservation.png}
    \caption{Screenshot of the preservation interface. Annotators are asked to read compare two passages and rate how similar the intent conveyed by the two passages is.}
    \label{fig:human-eval-preservation}
\end{figure*}

In this step, you will evaluate how much of the passage is attributable to one or more pieces of evidence (Figure \ref{fig:human-eval-attribution}).

In the interface, the passage of text and the context in which it was generated is shown on the left, and each piece of evidence is shown on the right. You will use all three (context, passage, evidence) to answer the following question for each sentence in the passage:
\textit{Is all of the information provided by this sentence fully supported by at least one piece of evidence?}

\paragraph{Determining the information provided by the sentence.}
Three points are key when determining information provided by the sentence:
\begin{enumerate}
\item The context and the other sentences of the passage are often critical in understanding the information provided by the sentence.
\item The context should only be used to understand the information provided by the sentence.
\item The evidence should be completely ignored for this step.
\end{enumerate}

Consider the following example:

\begin{quote}
\textit{Context:} who plays doug williams in days of our lives \\
\textit{Passage:} In the American daytime drama Days of Our Lives, Doug Williams and Julie Williams are portrayed by Bill Hayes and Susan Seaforth Hayes. \\
\end{quote}

In the above example, the meaning of the passage is clear even without seeing the query. But consider another example:

\begin{quote}
\textit{Context:} who plays doug williams in days of our lives \\
\textit{Passage:} he is played by Bill Hayes \\
\textit{Passage (interpreted):} \ul{Doug Williams} is played by Bill Hayes \ul{in days of our lives}
\end{quote}

In this case the pronoun ``he'' depends on the context, but it is clear that the  intended meaning of the passage can be reasonably interpreted as ``Doug Williams is played by Bill Hayes in days of our lives''. This interpretation is the ``information provided by the passage''. 

Pronouns such as he/she/it/they etc. are one case where context is needed to figure out the intended meaning of the system response. Here's another example (given with paraphrases of the information highlighted below):

\begin{quote}
\textit{Context:} when is the last time the us lost basketball at the olympics \\
\textit{Passage:} \ul{The last time they lost was in 2004}, when Argentina defeated the US 89–79. Most recently, they won gold in 2016. \\
\textit{Passage (interpreted):} The last time \ul{the United States lost basketball} at the Olympics was in 2004.
\end{quote}

The context should only be used to determine the information provided by the passage; at times, the passage may be about a slightly different topic than the context, for example:

\begin{quote}
\textit{Context:} the south west wind blows across nigeria between \\
\textit{Passage:} \ul{The Harmattan is a dry and dusty northeasterly trade wind that blows across West Africa} from December to March. It is very dusty because it blows across the Sahara.
\end{quote}

Here, the passage talks about a \textit{northeasterly wind}, while the context asks about a \textit{south-west wind}, but the passage can be fully understood.

In general, use your best judgment to determine the information provided by the passage. If the passage is hard to understand and you are unsure what the intended meaning of the passage is, \textit{mark the sentences as not attributed} and enter a comment with an explanation. As one example, take the following:

\begin{quote}
\textit{Context:} how many NBA championships did Michael Jordan win? \\
\textit{Passage:} it is the best team in the NBA
\end{quote}

\paragraph{Determining if the information accurately represents the evidence.}
Two points are key when determining whether the information accurately represents the evidence:
When interpreting a piece of evidence, use only the title and text of that specific evidence. Completely ignore the context, passage and all other evidence.
Check all the information in a sentence. If only some information is supported by the evidence, mark the sentence as not fully attributable.

Consider the following example:
\begin{quote}
\textit{Context:}
when did reba mcentire record back to god \\
\textit{Passage:}
\attrspan{Back to God was released by McEntire in 2017.} \\
\textit{Evidence:}
\ul{``Back to God'' is a song performed by American singer, Reba McEntire.} \ul{It was released} as the second single from her 2017 album, Sing it Now: Songs of Faith \& Hope, \ul{on January 20, 2017}.
\end{quote}
In the above example, it is reasonable to conclude that the evidence supports all the information in the passage, and we can mark the passage as being fully attributable. But consider another example:

\begin{quote}
\textit{Context:}
who won the womens 2017 ncaa basketball tournament \\
\textit{Passage:}
\nattrspan{South Carolina Gamecocks won the 2017 NCAA Women's Division I Basketball Tournament.} \\
\textit{Evidence:}
The \ul{South Carolina Gamecocks} defeated the Mississippi State Bulldogs, 67–55, to claim their first-ever \ul{national championship}.
\end{quote}

In this case, while the evidence also mentions the ``South Carolina Gamecocks'', it isn't clear that the national championship being mentioned is indeed the 2017 NCAA Women's Division I Basketball Tournament. The passage should be marked as not attributable.

Finally, when the passage contains multiple sentences, evaluate whether each sentence can be fully attributed to one or more pieces of evidence—it is possible for one sentence to be attributed while another is not. For example:

\begin{quote}
\textit{Context:}
who won the womens 2017 ncaa basketball tournament \\
\textit{Passage:}
\nattrspan{South Carolina Gamecocks won the 2017 NCAA Women's Division I Basketball Tournament.} \nattrspan{The final score is 67-55.} \attrspan{The championship game was held in Dallas, Texas.} \\
\textit{Evidence 1:}
The South Carolina Gamecocks defeated the Mississippi State Bulldogs, 67–55, to claim their first-ever national championship. \\
\textit{Evidence 2:}
The \ul{2017 NCAA Women's Division I Basketball Tournament} was played from Friday, March 17 to Sunday, April 2, 2017, with the Final Four played at the American Airlines Center in \ul{Dallas, Texas} on March 31 and April 2.
\end{quote}

The first two sentences cannot be attributed to either evidence for the same reason as the previous example, but the last sentence is fully supported by Evidence 2 and should be marked as attributed.

In general, you should use your best judgment in determining whether all of the information provided by the passage is ``an accurate representation of information in at least one evidence''. See Table \ref{tab:human-eval-attribution-examples} for additional examples.

We give the following final notes of guidance:
\begin{itemize}
\item  \textbf{Marking evidence as useful.} When reviewing each piece of evidence, mark it as useful if it helps you judge the attributability of any sentence, and mark it not useful if not. In the above example Evidence 1 is not useful because it didn't contain enough context to actually help you assess if the passage was attributable, but Evidence 2 was useful.
\item
\textbf{Contradicting evidence.} Mark a sentence as being attributed if any piece of evidence supports it: if two pieces of evidence contradict each other, but one of them supports the passage, mark the sentence as fully attributable.
\item
\textbf{More on the concept of ``accurate representation''.} We take as inspiration the journalist's conception of ``accurate representation''. For example, take \href{https://www.npr.org/about-npr/688139552/accuracy}{this excerpt on Accuracy in the NPR Ethics Handbook}: ``When quoting or paraphrasing anyone … consider whether the source would agree with the interpretation...'' In other words, if you had written the source document, consider whether you would view the system response as an accurate representation of information in that source document.
\end{itemize}


\begin{table*}
    \centering \small
    \begin{tabular}{@{}p{.25\textwidth} p{.45\textwidth} p{.25\textwidth}@{}}
    \toprule
        Context + Passage &
        Evidences &
        Notes \\ \midrule
        
\textit{Context: who played morticia in the addams family tv show}

\attrspan{The Addams Family is an American animated sitcom TV series.}
\nattrspan{It was first aired on NBC on September 24, 1973. Carolyn Jones played the role of Morticia.}
&
1/ The Addams Family (1973 TV series): The Addams Family is an American animated sitcom adaptation of the Charles Addams comic. The series was produced in 1973 and was rebroadcast the following season.

2/ The Addams Family (TV Series 1964–1966): When The Addams Family went off the air in 1966, network executives in charge of children's programming for NBC brought them back in 1973 for their own Saturday Morning cartoon show featuring the voices of Carolyn Jones from the original series.
&
While the evidence supports the show being aired in 1973, it doesn't specify the exact date (September 24, 1973).

Similarly, while the evidence mentions Carolyn Jones as being a voice actor, it doesn't say she played the role of Mortica.
\\      \midrule
\textit{
Context: when will the la sagrada familia be finished
}

\attrspan{The La Sagrada Familia is a large Roman Catholic church in Barcelona. It is designed by Antoni Gaudi.}
\nattrspan{It started construction in 1882, and the construction is still going on.}
\nattrspan{The estimated date to finish is 2026.}
&

1/ Sagrada Família - Wikipedia: The Basílica i Temple Expiatori de la Sagrada Família is a church in the Eixample district of Barcelona, Catalonia, Spain, and is currently the
largest unfinished Roman Catholic church.

2/ Find Out Sagrada Familia's Expected Finish Date: Visiting the breathtaking Sagrada Familia today also means witnessing the slow progress towards the completion of the project. Sagrada Familia is now  expected to be completed in 2026, the centenary of Gaudi's death.
It's a reasonable inference that La Sagrada Familia is the same as Sagrada Familia, even though the names differ slightly.
&
While Evidence 2 mentions Gaudi, it isn't clear this is a reference to Antoni Gaudi and further doesn't say that he designed the church.
\\
    \bottomrule
    \end{tabular}
    \caption{Additional examples for annotating attribution.}
    \label{tab:human-eval-attribution-examples}
\end{table*}

\subsection{Instructions: Intent Similarity}
In this step, you will evaluate how much similar the passage is to another passage (Figure \ref{fig:human-eval-preservation}).

In the interface, the passage A and passage B are both text generated by a system—given the same context in which it was generated. You will use all three (context, passage A, passage B) to answer the following question:
\textit{
How similar is the intent expressed by Passage A and Passage B? Please ignore any differences in details.
}

Two points are key when determining whether the two passages convey the same intent:
\begin{enumerate}
    \item  Judge the similarity solely based on the similarity in the type and quantity of information provided by each passage.
    \item Ignore any differences in factual details between the two passages.
\end{enumerate}

Consider the following examples:

\begin{quote}
\textit{Context:}
who pays medical bills in great britain where does the money come from to pay these bills \\
\textit{Passage A:}
Britain's National Health Service (NHS) is paid for through general taxation \ul{and national insurance}. In 2017/18, the NHS budget was £\ul{176.5bn}. \\
\textit{Passage B:}
Britain's National Health Service (NHS) is paid for through general taxation. In 2017/18, the NHS budget was £\ul{118bn}. \\
\textit{Rating:} \ul{Very similar.} Passage A is about the same topic as Passage B, with a similar level of detail and style of presentation. They may differ in factual details.
\end{quote}

The above example should be rated ``very similar'' because both passages include information about (1) how the NHS is paid for, and (2) what its budget in 2017/18 was, though they differ in their actual answers to these questions.

\begin{quote}
\textit{Context:}
who is the owner of reading football club \\
\textit{Passage A:}
Reading's owner is \ul{Yongge Dai}. Yongge \ul{Dai is also the president of Chinese company Dai Yongge Real Estate. Yongge's son, Dai Xiu Li, is Reading's vice-president.} \\
\textit{Passage B:}
Reading's owner is \ul{Dai Yongge}. Yongge's \ul{brother and sister pair behind the Reading FC takeover---Dai Yongge and Dai Xiu Li---has made their fortune through a massive property empire. Mr Dai, has been the chairman of Renhe Commercial since 1999, which is an organisation owned by his sister behind a vast network of underground shopping centres in China}. \\
\textit{Rating:}
\ul{Somewhat similar.} Passage A is about the same topic as Passage B, but differs substantially in level of detail or style of presentation. They may differ in factual details.
\end{quote}

The above example should be rated ``somewhat similar'' because both passages are still about the same topic—Reading's owner— but differ substantially in the information they discuss: Passage A includes information about (1a) who Reading's owner is, (2a) which company they are the president of and (3a) who their vice-president is. In contrast, while Passage B shares information about (1a), it also includes information about (2b) how the Reading owner made their fortune, (3b) their company position and how long they held it for and (4b) what the company also owns.


\begin{quote}
\textit{Context:}
what is the numbers of total elected member of indian parliment in present time \\
\textit{Passage A:}
The total number of elected members of the \ul{Lok} Sabha is \ul{543}. \\
\textit{Passage B:}
The total number of elected members of the \ul{Rajya} Sabha is \ul{238}. \\
\textit{Rating:}
\ul{Not at all similar.} Passage A is about a significantly different topic than Passage B.
\end{quote}

Even though the passages look very similar, the above example should be rated ``not at all similar'' because the two passages are about significantly different topics: ``the Lok Sabha'' vs ``the Rajya Sabha''.
